% Курсова: один документ (титул + пояснювальна записка). Оформлення те саме, що в дипломі
% (ДСТУ 3008:2015); склад розділів звірити з методичкою курсової своєї дисципліни.
\documentclass[coursework]{kpi-thesis}
\usepackage{kpi-titles}
\kpiset{
  title         = {\todo{Тема курсової роботи}},
  discipline    = {\todo{Назва дисципліни}},
  group         = {ІП-XX},
  course        = {IV},
  gender        = {m},
  student       = {\todo{Прізвище Ім'я По батькові}},
  advisor       = {\todo{посада, науковий ступінь, Прізвище І. П.}},
  year          = {2027},
}

\begin{document}
\kpiCourseworkTitle
\mainmatter
\tableofcontents

\structural{ВСТУП}

\textbf{Актуальність.} \todo{...} \cite{example}.

\textbf{Мета роботи} -- \todo{практичний результат}.

Для досягнення мети потрібно:
\begin{itemize}
\item \todo{завдання};
\end{itemize}

\section{Аналіз предметної області та постановка задачі}

\todo{...}

\subsection{Висновки до розділу}
\todo{...}

\section{Проєктування програмного забезпечення}

\todo{...}

\subsection{Висновки до розділу}
\todo{...}

\section{Реалізація програмного забезпечення}

\todo{...}

\subsection{Висновки до розділу}
\todo{...}

\section{Тестування програмного забезпечення}

\todo{...}

\subsection{Висновки до розділу}
\todo{...}

\structural{ВИСНОВКИ}

\todo{...}

% ДСТУ 8302:2015, за порядком першої згадки. Зразки — skills/kpi-thesis/references/requirements.md.
\begin{thebibliography}{99}
\bibitem{example} \todo{Автор А. Б. Назва. Місто : Видавництво, 2025. 100 с.}
\end{thebibliography}


\appendices
\appendix{Лістинг програми}
% \lstinputlisting{../src/main.py}
\todo{лістинг програми}
\end{document}
